% ==============================================================================
% Course Syllabus: MSQF 532 Financial Mathematics
% Formatted to fit on a single, clean official syllabus page
% ==============================================================================

\newpage
\thispagestyle{plain}

\begin{center}
    {\large\textbf{MSQF 532: FINANCIAL MATHEMATICS}} \hfill {\large\textbf{CREDITS: 4}}
\end{center}
\vspace{-0.6em}
\hrule height 0.8pt
\vspace{0.6em}

\noindent\textbf{Course Objectives:}
\begin{itemize}[leftmargin=1.5em, label=$\bullet$, noitemsep, topsep=1pt, parsep=1pt]
    \item \textit{To equip students with a solid understanding of fundamental mathematical concepts and techniques used in financial analysis, valuation of financial instruments, and portfolio management.}
    \item \textit{To develop the ability to apply quantitative methods for solving practical problems in finance, including risk assessment, derivatives pricing, and financial decision-making.}
\end{itemize}

\vspace{0.4em}
\noindent\textbf{Course Outcomes (COs):}
\begin{description}[leftmargin=1.4cm, labelsep=0.4em, itemsep=1pt, topsep=1pt, parsep=0pt, font=\normalfont\textbf]
    \item[CO1:] Apply financial mathematics to evaluate investments and corporate finance decisions
    \item[CO2:] Use quantitative tools for bond, equity, and portfolio valuation
    \item[CO3:] Analyze risk and return to support financial strategy
    \item[CO4:] Interpret and construct financial statements for cash flow management
    \item[CO5:] Apply mathematical models to real-world financial problems
\end{description}

\vspace{0.4em}
\noindent\textbf{Unit I: Introduction to Financial Mathematics}\\
\small
Role and importance of mathematics in finance --- Time value of money concepts --- Simple interest and compound interest --- Nominal and effective rates --- Discounting and present value --- Basic annuities: ordinary annuities, annuities due --- Growth and decay curves.
\normalsize

\vspace{0.3em}
\noindent\textbf{Unit II: Valuation of Financial Instruments}\\
\small
Present value and future value of annuities --- Perpetuities and growing annuities --- Loan amortization and sinking funds --- Bond valuation: yield to maturity (YTM), yield to call (YTC) --- Equity valuation: dividend discount model (DDDM), price-earnings model.
\normalsize

\vspace{0.3em}
\noindent\textbf{Unit III: Financial Ratios and Cash Flow Analysis}\\
\small
Key financial ratios: liquidity, profitability, leverage, efficiency --- Application of ratios in decision making --- Cash flow statements: operating, investing, and financing cash flows --- Free cash flow analysis --- Break-even and sensitivity analysis.
\normalsize

\vspace{0.3em}
\noindent\textbf{Unit IV: Risk, Return, and Portfolio Mathematics}\\
\small
Measurement of return: arithmetic vs geometric mean --- Standard deviation and variance of returns --- Coefficient of variation --- Covariance and correlation --- Portfolio return and risk --- Diversification and optimal portfolio construction --- Capital Asset Pricing Model (CAPM).
\normalsize

\vspace{0.3em}
\noindent\textbf{Unit V: Advanced Topics and Applications}\\
\small
Linear programming in financial decision-making --- Introduction to financial derivatives (options, forwards, futures) --- Option valuation basics: Binomial and Black-Scholes models --- Duration and convexity of bonds --- Probability distributions in finance: normal, log-normal --- Simulation and scenario analysis in financial planning.
\normalsize

\vspace{0.5em}
\noindent\textbf{Books for Reference:}
\begin{enumerate}[leftmargin=1.5em, label=\arabic*., itemsep=1pt, topsep=1pt, parsep=0pt]
    \small
    \item Bodie, Z., Kane, A., \& Marcus, A. J. (2022). \textit{Investments} (12th ed.). McGraw-Hill Education.
    \item Tuckman, B., \& Serrat, A. (2023). \textit{Fixed Income Securities: Tools for Today's Markets} (4th ed.). Wiley.
    \item Fabozzi, F. J. (2023). \textit{Financial Mathematics: A Comprehensive Treatment} (2nd ed.). Wiley.
    \item Hull, J. C. (2024). \textit{Options, Futures, and Other Derivatives} (11th ed.). Pearson.
    \item Ross, S. A., Westerfield, R. W., Jaffe, J., \& Jordan, B. D. (2022). \textit{Corporate Finance} (13th ed.). McGraw-Hill Education.
\end{enumerate}

\vspace{0.4em}
\begin{center}
\footnotesize
\begin{tabular}{|c|c|c|c|c|c|}
\hline
\textbf{COs} & \textbf{PO1} & \textbf{PO2} & \textbf{PO3} & \textbf{PO4} & \textbf{PO5} \\
\hline
\textbf{CO1} & XXX & XXX & XXX & XX & XX \\
\hline
\textbf{CO2} & XXX & XXX & XXX & XX & XX \\
\hline
\textbf{CO3} & XXX & XX & XXX & XXX & XX \\
\hline
\textbf{CO4} & XX & XX & XX & XXX & XX \\
\hline
\textbf{CO5} & XXX & XXX & XXX & XX & XXX \\
\hline
\end{tabular}
\end{center}
\normalsize

\newpage
